\documentclass[11pt]{article}

% write-latex compliant: no fontspec, no XeLaTeX-only packages.
\usepackage[utf8]{inputenc}
\usepackage[T1]{fontenc}
\usepackage{lmodern}
\usepackage[margin=1in]{geometry}
\usepackage{amsmath, amssymb, amsthm, mathtools}
\usepackage{hyperref}
\hypersetup{colorlinks=true, linkcolor=blue, urlcolor=blue, citecolor=blue}

\newtheorem{theorem}{Theorem}
\newtheorem{lemma}[theorem]{Lemma}
\newtheorem{corollary}[theorem]{Corollary}
\theoremstyle{definition}
\newtheorem{definition}[theorem]{Definition}
\theoremstyle{remark}
\newtheorem{remark}[theorem]{Remark}

\title{Control variates cancel the leading client-heterogeneity drift in a federated two-step round}
\author{Companion proof for: \href{https://github.com/pleyva2004/communication-efficient-learning-of-deep-networks}{pleyva2004/communication-efficient-learning-of-deep-networks} \\ \small Improvement slug: \texttt{heterogeneity-control-variate}}
\date{}

\begin{document}
\maketitle

\iffalse
NOTE: this file is a per-improvement proof, generated by the study-paper skill, Stage 7.
PROOF mode is selected when the improvement is mathematical, theoretical, or a
complexity-class statement (see Stage 7 heuristic). The proof must be complete:
no "obvious", no "left as exercise", no "by standard arguments".
This file loads NO natbib/bibliography; all prior work is referenced as inline plaintext.
\fi

\begin{abstract}
The math deep dive (eq.~(3.5)) showed that running two local gradient-descent steps per
client and then averaging deviates from two genuine centralized gradient-descent steps by a
leading term $\eta^2\,\mathrm{Cov}_k(H_k,g_k)+O(\eta^3)$, the client-weighted covariance
between local Hessians $H_k$ and local gradients $g_k$. This is the ``client-drift'' object
that later federated-learning theory (e.g.\ Karimireddy et al.\ 2020, SCAFFOLD) is explicitly
designed to remove. Here we prove, by carrying the \emph{same} Taylor expansion but replacing
each local gradient $g_{\mathrm{local}}(w)$ with the control-variate-corrected direction
$g_{\mathrm{local}}(w)-c_k+c$ (where $c_k$ is a fixed estimate of client $k$'s gradient at the
shared point $w_t$ and $c=\sum_k\alpha_k c_k$ is its client-weighted mean), that the entire
order-$\eta^2$ covariance term cancels. Consequently the corrected averaged two-step round
matches two centralized gradient-descent steps up to $O(\eta^3)$, recovering, to leading
order, the per-round progress of centralized full-batch gradient descent even on an
arbitrarily heterogeneous (non-IID) partition.
\end{abstract}

\section*{Setup and notation}

We adopt verbatim the setup of the math deep dive. The global objective is the finite sum
$f(w)=\sum_{k=1}^{K}\alpha_k F_k(w)$ with mixture weights $\alpha_k:=n_k/n\ge0$,
$\sum_{k=1}^{K}\alpha_k=1$, where $F_k(w)=\frac{1}{n_k}\sum_{i\in P_k}f_i(w)$ is client $k$'s
local objective over its example set $P_k$ of size $n_k$; this is the exact algebraic
identity (1.1) of the deep dive, valid for any partition. All clients begin a round at the
common shared point $w_t\in\mathbb{R}^d$ (the load-bearing shared-initialization assumption).
We fix the round and write everything relative to $w_t$.

We work in the \emph{quadratic / second-order local model} used to derive eq.~(3.5): each
$F_k$ is twice continuously differentiable in a neighborhood of $w_t$, and we expand each
local gradient to first order about $w_t$ with a uniform $O(\cdot)$ remainder. Concretely,
writing $g_k:=\nabla F_k(w_t)$ and $H_k:=\nabla^2 F_k(w_t)$, for any displacement $\delta$ of
size $O(\eta)$ Taylor's theorem with remainder gives
\begin{equation}
\nabla F_k(w_t+\delta)=g_k+H_k\,\delta+R_k(\delta),\qquad \|R_k(\delta)\|=O(\|\delta\|^2),
\label{eq:taylor}
\end{equation}
where the constant hidden in $O(\|\delta\|^2)$ is uniform over the finitely many clients
$k\in\{1,\dots,K\}$ (take the maximum over $k$ of the per-client constants). Because every
displacement we encounter is $O(\eta)$, every such remainder contributes at order $\eta^2$
\emph{inside} a term that is already multiplied by $\eta$, i.e.\ at order $\eta^3$ in the
position iterate; this is exactly the bookkeeping that makes the deep dive's $O(\eta^3)$
statement precise. The learning rate is $\eta>0$, taken small enough that the round's
displacements stay in the neighborhood where \eqref{eq:taylor} holds.

\begin{definition}[Working notation]
Throughout the round:
\begin{itemize}
  \item $g_k:=\nabla F_k(w_t)\in\mathbb{R}^d$ --- client $k$'s local gradient at the shared
        point $w_t$.
  \item $H_k:=\nabla^2 F_k(w_t)\in\mathbb{R}^{d\times d}$ --- client $k$'s local Hessian at
        $w_t$.
  \item $\nabla f:=\nabla f(w_t)=\sum_{k=1}^{K}\alpha_k g_k$ --- the global gradient at $w_t$;
        the equality is the exact gradient-aggregation identity (2.2) of the deep dive,
        obtained by differentiating $f=\sum_k\alpha_k F_k$.
  \item $H:=\nabla^2 f(w_t)=\sum_{k=1}^{K}\alpha_k H_k$ --- the global Hessian at $w_t$, again
        by differentiating the mixture identity twice.
  \item $\mathrm{Cov}_k(H_k,g_k):=\sum_{k=1}^{K}\alpha_k H_k g_k-\Big(\sum_{k=1}^{K}\alpha_k
        H_k\Big)\Big(\sum_{k=1}^{K}\alpha_k g_k\Big)=\sum_{k=1}^{K}\alpha_k H_k g_k-H\,\nabla f$
        --- the client-weighted covariance between local Hessians and local gradients (a
        vector in $\mathbb{R}^d$), with the $\alpha_k$ acting as the client distribution. This
        is the leading deviation term of eq.~(3.5).
  \item \textbf{Control variates.} Each client $k$ carries a fixed vector $c_k$ that estimates
        its own local gradient direction at the shared point $w_t$; in the static, exact-estimate
        regime analyzed here we take $c_k:=g_k=\nabla F_k(w_t)$, held \emph{constant} for the
        whole round (it is computed once at $w_t$ and not refreshed between the two local
        steps). The server control variate is the client-weighted mean
        $c:=\sum_{k=1}^{K}\alpha_k c_k=\sum_{k=1}^{K}\alpha_k g_k=\nabla f$, also held constant
        for the round.
  \item \textbf{Corrected local direction.} At any iterate $w$, instead of stepping along the
        raw local gradient $g_{\mathrm{local}}(w):=\nabla F_k(w)$, client $k$ steps along the
        control-variate-corrected direction
        \begin{equation}
          \tilde g_k(w):=\nabla F_k(w)-c_k+c.
          \label{eq:corrected-dir}
        \end{equation}
\end{itemize}
\end{definition}

\section*{Theorem}

\begin{theorem}[Control variates cancel the order-$\eta^2$ heterogeneity drift]
\label{thm:main}
Fix a round at the shared point $w_t$ and adopt the quadratic / second-order local model
\eqref{eq:taylor} with static, exact control variates $c_k=g_k$ and
$c=\sum_k\alpha_k g_k=\nabla f$. Let each client run \emph{two} local steps using the corrected
direction \eqref{eq:corrected-dir},
\[
  w^{k}_{(0)}=w_t,\qquad
  w^{k}_{(1)}=w^{k}_{(0)}-\eta\,\tilde g_k\!\big(w^{k}_{(0)}\big),\qquad
  w^{k}_{(2)}=w^{k}_{(1)}-\eta\,\tilde g_k\!\big(w^{k}_{(1)}\big),
\]
and let the server form the client-weighted average $\bar w_{t+2}:=\sum_{k=1}^{K}\alpha_k
w^{k}_{(2)}$. Let $G^{(2)}(w_t)$ denote two genuine centralized gradient-descent steps on $f$
from $w_t$, namely $u_1=w_t-\eta\,\nabla f$, $G^{(2)}(w_t)=u_1-\eta\,\nabla f(u_1)$. Then the
order-$\eta^2$ deviation $\eta^2\,\mathrm{Cov}_k(H_k,g_k)$ of eq.~(3.5) vanishes to first
order: the corrected averaged round equals two centralized steps up to $O(\eta^3)$,
\[
  \bar w_{t+2}
  = w_t-2\eta\,\nabla f+\eta^2 H\,\nabla f+O(\eta^3)
  = G^{(2)}(w_t)+O(\eta^3),
\]
whereas the \emph{uncorrected} FedAvg round (using $g_{\mathrm{local}}$ directly) satisfies
$\bar w^{\mathrm{FedAvg}}_{t+2}=G^{(2)}(w_t)+\eta^2\,\mathrm{Cov}_k(H_k,g_k)+O(\eta^3)$. In
particular the leading discrepancy is reduced from $\Theta(\eta^2)$ (generically nonzero under
client heterogeneity) to $O(\eta^3)$.
\end{theorem}

\section*{Proof}

\begin{proof}
The proof carries the same Taylor expansion as the deep dive's derivation of eq.~(3.5), but
applied to the corrected direction \eqref{eq:corrected-dir}. We proceed in five steps: (i)
compute the corrected first step; (ii) Taylor-expand the corrected second step; (iii) assemble
the corrected two-step endpoint per client; (iv) average over clients; (v) compare against two
centralized steps and against the uncorrected FedAvg round, exhibiting the cancellation of the
covariance term.

\medskip
\noindent\textbf{Step (i): the corrected first step is identical across clients and equals one
centralized step.}
Evaluate \eqref{eq:corrected-dir} at the shared start $w^{k}_{(0)}=w_t$. By definition
$\nabla F_k(w_t)=g_k$ and, in the static exact-estimate regime, $c_k=g_k$, so the raw local
gradient and the client's own control variate cancel exactly:
\begin{equation}
  \tilde g_k(w_t)=\underbrace{\nabla F_k(w_t)}_{=\,g_k}-\underbrace{c_k}_{=\,g_k}+\,c
  = g_k-g_k+c = c = \nabla f.
  \label{eq:first-dir}
\end{equation}
This is the crux of the correction: the term $-c_k$ removes the client's own gradient and
$+c$ substitutes the \emph{global} direction $\nabla f$, so the first corrected step is the
same for every client and equals exactly one centralized gradient-descent step,
\begin{equation}
  w^{k}_{(1)}=w_t-\eta\,\tilde g_k(w_t)=w_t-\eta\,\nabla f
  \qquad\text{for every }k.
  \label{eq:w1}
\end{equation}
In particular $w^{k}_{(1)}$ does not depend on $k$; we may write it as $u_1:=w_t-\eta\,\nabla
f$, which is precisely the first centralized iterate.

\medskip
\noindent\textbf{Step (ii): Taylor-expand the corrected second direction.}
We now evaluate \eqref{eq:corrected-dir} at $w^{k}_{(1)}=w_t-\eta\,\nabla f$. The displacement
from $w_t$ is $\delta:=-\eta\,\nabla f$, which is $O(\eta)$, so \eqref{eq:taylor} applies with
remainder $\|R_k(\delta)\|=O(\eta^2)$:
\begin{equation}
  \nabla F_k\big(w_t-\eta\,\nabla f\big)
  = g_k + H_k(-\eta\,\nabla f) + R_k(-\eta\,\nabla f)
  = g_k - \eta\,H_k\,\nabla f + O(\eta^2).
  \label{eq:grad-at-w1}
\end{equation}
Subtract the static control variate $c_k=g_k$ and add $c=\nabla f$. The constant local
gradient $g_k$ cancels against $-c_k=-g_k$ a \emph{second} time --- this is the structural
reason the correction keeps working past the first step --- leaving
\begin{equation}
  \tilde g_k\big(w^{k}_{(1)}\big)
  = \big(g_k-\eta\,H_k\,\nabla f+O(\eta^2)\big) - g_k + \nabla f
  = \nabla f - \eta\,H_k\,\nabla f + O(\eta^2).
  \label{eq:second-dir}
\end{equation}
Notice that the leading $O(1)$ part of the corrected second direction is again the global
$\nabla f$, with the \emph{only} client-dependent piece appearing at order $\eta$ through the
client's own Hessian $H_k$ acting on the \emph{global} gradient $\nabla f$ --- not on the
local $g_k$. This single change (the matrix $H_k$ now multiplies $\nabla f$ rather than $g_k$)
is what will make the covariance collapse in Step (v).

\medskip
\noindent\textbf{Step (iii): assemble the corrected two-step endpoint per client.}
Combine \eqref{eq:w1} and \eqref{eq:second-dir}:
\begin{align}
  w^{k}_{(2)}
  &= w^{k}_{(1)} - \eta\,\tilde g_k\big(w^{k}_{(1)}\big) \notag\\
  &= \big(w_t-\eta\,\nabla f\big) - \eta\big(\nabla f-\eta\,H_k\,\nabla f+O(\eta^2)\big) \notag\\
  &= w_t - \eta\,\nabla f - \eta\,\nabla f + \eta^2 H_k\,\nabla f + \eta\cdot O(\eta^2) \notag\\
  &= w_t - 2\eta\,\nabla f + \eta^2 H_k\,\nabla f + O(\eta^3).
  \label{eq:w2-client}
\end{align}
Here the term $\eta\cdot O(\eta^2)=O(\eta^3)$ absorbs the propagated Taylor remainder: the
$O(\eta^2)$ error in the second direction \eqref{eq:second-dir} is multiplied by the step size
$\eta$, hence enters the position iterate \eqref{eq:w2-client} only at order $\eta^3$, exactly
as claimed. This is the corrected analogue of the deep dive's uncorrected endpoint
$w^{k}_{(2),\mathrm{FedAvg}}=w_t-2\eta g_k+\eta^2 H_k g_k+O(\eta^3)$: the $O(1)$ term has moved
from the client-specific $-2\eta g_k$ to the global $-2\eta\,\nabla f$, and the $O(\eta^2)$
term from $H_k g_k$ to $H_k\,\nabla f$.

\medskip
\noindent\textbf{Step (iv): average over clients.}
Apply the server's client-weighted average $\sum_k\alpha_k(\cdot)$ to \eqref{eq:w2-client}.
Linearity of the sum and $\sum_k\alpha_k=1$ give, term by term,
\begin{align}
  \bar w_{t+2}
  &= \sum_{k=1}^{K}\alpha_k\Big(w_t-2\eta\,\nabla f+\eta^2 H_k\,\nabla f+O(\eta^3)\Big) \notag\\
  &= \Big(\underbrace{\textstyle\sum_k\alpha_k}_{=1}\Big)w_t
     - 2\eta\Big(\underbrace{\textstyle\sum_k\alpha_k}_{=1}\Big)\nabla f
     + \eta^2\Big(\underbrace{\textstyle\sum_k\alpha_k H_k}_{=\,H}\Big)\nabla f
     + O(\eta^3) \notag\\
  &= w_t - 2\eta\,\nabla f + \eta^2 H\,\nabla f + O(\eta^3).
  \label{eq:avg}
\end{align}
Two facts make the averaging clean. First, $w_t$ and $\nabla f$ are constants with respect to
the index $k$ (the global gradient $\nabla f=\sum_j\alpha_j g_j$ does not depend on the outer
$k$), so they pass through the average unchanged via $\sum_k\alpha_k=1$. Second, the only
$k$-dependent object surviving to order $\eta^2$ is the Hessian $H_k$, and it is multiplied by
the \emph{$k$-independent} vector $\nabla f$; hence the average of $H_k\,\nabla f$ is exactly
$\big(\sum_k\alpha_k H_k\big)\nabla f=H\,\nabla f$, with no cross term. The $O(\eta^3)$
remainders, being a finite $\alpha_k$-weighted ($\sum_k\alpha_k=1$, $\alpha_k\ge0$) average of
$O(\eta^3)$ vectors with uniform constant, remain $O(\eta^3)$.

\medskip
\noindent\textbf{Step (v): compare to two centralized steps, and exhibit the cancellation.}
Two genuine centralized gradient-descent steps on $f$ from $w_t$ give, by the same Taylor
expansion \eqref{eq:taylor} applied to $\nabla f$ with global Hessian $H=\nabla^2 f(w_t)$,
\begin{align}
  u_1&=w_t-\eta\,\nabla f,\notag\\
  G^{(2)}(w_t)&=u_1-\eta\,\nabla f(u_1)
   = u_1-\eta\big(\nabla f-\eta H\,\nabla f+O(\eta^2)\big)\notag\\
  &= w_t-2\eta\,\nabla f+\eta^2 H\,\nabla f+O(\eta^3),
  \label{eq:centralized}
\end{align}
which is identical, term by term up to $O(\eta^3)$, to the corrected averaged round
\eqref{eq:avg}. Subtracting,
\begin{equation}
  \bar w_{t+2}-G^{(2)}(w_t)=O(\eta^3).
  \label{eq:corrected-gap}
\end{equation}
To see precisely what was cancelled, repeat Step (iv) for the \emph{uncorrected} FedAvg
endpoint $w^{k}_{(2),\mathrm{FedAvg}}=w_t-2\eta g_k+\eta^2 H_k g_k+O(\eta^3)$ of the deep dive.
Averaging,
\begin{equation}
  \bar w^{\mathrm{FedAvg}}_{t+2}
  = w_t-2\eta\!\sum_k\alpha_k g_k+\eta^2\!\sum_k\alpha_k H_k g_k+O(\eta^3)
  = w_t-2\eta\,\nabla f+\eta^2\!\sum_k\alpha_k H_k g_k+O(\eta^3),
  \label{eq:fedavg-avg}
\end{equation}
using $\sum_k\alpha_k g_k=\nabla f$. Subtracting the centralized expansion
\eqref{eq:centralized} from \eqref{eq:fedavg-avg}, the $w_t$ and $-2\eta\,\nabla f$ terms cancel
and the leading deviation is
\begin{equation}
  \bar w^{\mathrm{FedAvg}}_{t+2}-G^{(2)}(w_t)
  = \eta^2\Big(\sum_{k=1}^{K}\alpha_k H_k g_k-H\,\nabla f\Big)+O(\eta^3)
  = \eta^2\,\mathrm{Cov}_k(H_k,g_k)+O(\eta^3),
  \label{eq:fedavg-gap}
\end{equation}
which is exactly eq.~(3.5) of the deep dive, by the definition of $\mathrm{Cov}_k$ recorded in
the Working notation.

The cancellation is now transparent at the level of the order-$\eta^2$ coefficient. The
uncorrected coefficient is $\sum_k\alpha_k H_k g_k$, in which the matrix $H_k$ multiplies the
\emph{client-specific} gradient $g_k$; its average does \emph{not} factor, leaving the residual
$\sum_k\alpha_k H_k g_k-H\,\nabla f=\mathrm{Cov}_k(H_k,g_k)$, which is generically nonzero
whenever the clients are heterogeneous (different $H_k$ correlated with different $g_k$). The
corrected coefficient, by contrast, is $\sum_k\alpha_k H_k\,\nabla f$, in which $H_k$
multiplies the \emph{global} gradient $\nabla f$ --- a constant in $k$ --- so the average
factors exactly as $\big(\sum_k\alpha_k H_k\big)\nabla f=H\,\nabla f$, leaving no residual.
Equivalently, the difference between the corrected and uncorrected order-$\eta^2$ terms is
\begin{equation}
  \eta^2\sum_{k=1}^{K}\alpha_k H_k\big(\nabla f-g_k\big)
  = \eta^2\Big(H\,\nabla f-\sum_{k=1}^{K}\alpha_k H_k g_k\Big)
  = -\,\eta^2\,\mathrm{Cov}_k(H_k,g_k),
  \label{eq:explicit-removal}
\end{equation}
i.e.\ the control variate replaces the local gradient $g_k$ that $H_k$ would have acted on by
the global gradient $\nabla f$ inside the second step, and the resulting change is precisely
$-\,\eta^2\,\mathrm{Cov}_k(H_k,g_k)$, the exact negative of the FedAvg drift
\eqref{eq:fedavg-gap}. Adding this to \eqref{eq:fedavg-gap} sends the leading deviation to
zero, leaving only $O(\eta^3)$, which is \eqref{eq:corrected-gap}.

This proves both displayed equalities of Theorem~\ref{thm:main}: the corrected averaged round
matches two centralized steps to $O(\eta^3)$, while the uncorrected round retains the
$\Theta(\eta^2)$ heterogeneity term. The order-$\eta^2$ deviation
$\eta^2\,\mathrm{Cov}_k(H_k,g_k)$ therefore vanishes to first order under the control-variate
correction.
\end{proof}

\section*{Discussion}

\paragraph{What this gives us.}
The leading-order client-heterogeneity penalty of running more than one local step --- the term
$\eta^2\,\mathrm{Cov}_k(H_k,g_k)$ that the deep dive identified as the entire $O(\eta^2)$ price
of the second local step on non-IID data --- is removed exactly by the control-variate
correction $g_{\mathrm{local}}(w)\mapsto g_{\mathrm{local}}(w)-c_k+c$. After the fix, a
corrected two-step federated round reproduces two centralized full-batch gradient-descent steps
up to $O(\eta^3)$, i.e.\ the per-round progress is, to leading order, indistinguishable from
centralized training \emph{regardless of how heterogeneous the partition is}. The mechanism is
explicit and local: $-c_k$ subtracts the client's own gradient and $+c$ inserts the global one,
so each corrected step is centered on $\nabla f$ rather than on the client's biased $g_k$. This
is the formal content behind importing SCAFFOLD-style control variates (Karimireddy et al.\
2020) into FedAvg, and it is the leading-order counterpart of the proximal anchor used by
FedProx (Li et al.\ 2020), which attacks the same drift object by a different route.

\paragraph{What this does not give us.}
The result is a \emph{leading-order}, single-round, per-step statement, not a convergence
theorem: it does not by itself bound the number of communication rounds, prove descent of $f$,
or rule out divergence; it only equates one corrected two-step round with two centralized steps
to $O(\eta^3)$. It says nothing about the $O(\eta^3)$ and higher terms, which still carry
genuine client-dependence (third derivatives of $F_k$, the propagated remainder
$\eta\cdot O(\eta^2)$ of Step (iii)); the correction zeroes the $\eta^2$ coefficient, not the
whole gap. The analysis is deterministic and full-batch, with $g_{\mathrm{local}}=\nabla F_k$
exact; it does not address stochastic minibatch noise (where $c_k$, $c$ are themselves noisy
estimates and the cancellation holds only in expectation), partial client participation, or the
extra per-client state and the control-variate refresh that a multi-round deployment requires.
The proof also assumes the standard FedAvg shared initialization $w^{k}_{(0)}=w_t$; without it
even the first step \eqref{eq:w1} would differ across clients.

\paragraph{Where it would break.}
Two assumptions are load-bearing. (1) \emph{Static, accurate control variates.} The
cancellation in Steps (i)--(ii) relied twice on $c_k=g_k=\nabla F_k(w_t)$ holding for the whole
round. If $c_k$ is \emph{stale} --- estimated at an earlier global model $w_{t'}\neq w_t$, as
happens with intermittent participation --- then $c_k=\nabla F_k(w_{t'})=g_k-H_k(w_t-w_{t'})+
\dots$, and the leftover $-c_k+c$ no longer cancels $g_k$ cleanly: a residual proportional to
the staleness gap $w_t-w_{t'}$ re-enters at order $\eta$ in the corrected direction and hence at
order $\eta^2$ in the iterate, partially restoring a drift term. The cleaner $c_k=g_k$ the
estimate, the more of the $\eta^2$ covariance is removed; an arbitrarily stale $c_k$ recovers
no benefit. (2) \emph{Mild nonlinearity (validity of the quadratic model).} The whole argument
truncates the Taylor expansion \eqref{eq:taylor} at first order in the gradient (second order in
$F_k$). Under strong nonlinearity --- large third derivatives $\nabla^3 F_k$, or a learning rate
$\eta$ large enough that the round's displacements leave the neighborhood where
\eqref{eq:taylor} holds --- the discarded $O(\eta^3)$ terms are no longer negligible and carry
their own uncancelled client-heterogeneity contributions; the control variate is then guaranteed
to remove only the $\eta^2$ piece, and the practical benefit shrinks as $\eta$ grows or the
local landscapes become more curved/heterogeneous at third order. Pushing the number of local
steps well beyond two compounds this: each additional step adds higher-order remainders that the
single first-order control variate does not address.

\paragraph{Empirical companion.}
This proof is the validation file for the improvement \texttt{heterogeneity-control-variate} in
\texttt{05-improvements.tex} (subsection M.1). If a measurement-mode prototype also exists at
\texttt{improvements/heterogeneity-control-variate.py}, it serves as a sanity check, not as
evidence --- the theorem above is the load-bearing claim.

\end{document}
